\subsection{Законы больших чисел.}
Пусть дана некоторая последовательность случайных величин $\{\xi_n\}$.
Будем обозначать за $\overline{\xi}_n$ следующую случайную величину:
\[
    \overline{\xi_n} = \dfrac{1}{n}\sum\limits_{k = 1}^{n}\xi_k.
\]
Далее мы будем говорить что $\xi_n$ подчиняется ЗБЧ, если существует числовая последовательность $\{ a_n \}$:
\[
    \overline{\xi}_n - a_n \xrightarrow{\P} 0, n \ra \infty.
\]
Здесь и далее будем работать с $a_n = \Ex \overline{\xi_n} = \dfrac{1}{n}\sum\limits_{k = 1}^n \Ex \xi_k$.
С такой последовательностью $a_n$ среднее $\xi_n$ примет вид
\[
    \overline{\xi_n} - a_n = \overline{\xi_n} - \dfrac{1}{n}\sum\limits_{k = 1}^{n} \Ex \xi_k = \overline{\overset{\circ}{\xi_n}}.
\]
По сути, глобальной задачей является описание необходимых и достаточных условий для выполнения ЗБЧ.
\subsubsection{Достаточные условия выполнення ЗБЧ в форме Маркова, Чебышёва и Хинчина.}
\begin{theorem}[ЗБЧ в форме Маркова]
    Если выполнено следующее условие на дисперсии $\xi_n$, то выполнен ЗБЧ:
    \[
        \dfrac{1}{n^2}\D\biggr(\sum\limits_{k = 1}^{n} \xi_k\biggr) \ra 0, n \ra +\infty.
    \]
\end{theorem}
\begin{proof}
    Пусть $\epsilon > 0$.
    Используем неравенство Чебышёва для последовательности $\overline{{\xi_n}}$ и получим:
    \[
        \P(|\overline{\xi_n} - \Ex \overline{\xi_n}| \geq \epsilon) \leq \dfrac{1}{n^2\epsilon^2}{\D\left( \overline{\xi_n}\right)} = \dfrac{1}{n^2\epsilon^2}{\D\biggr( \sum\limits_{k = 1}^{n} \xi_k \biggr)} \ra 0, n \ra +\infty.
    \]
    Что и означает выполнение ЗБЧ.
\end{proof}
\begin{corollary}[ЗБЧ в форме Чебышёва]
    Пусть $\{\xi_n\}$~---~некоррелированные случайные величины, также $\exists \alpha \in (0, 1)$ и $C > 0$: $\D \xi_n \leq C n^{\alpha}$. Тогда выполнен ЗБЧ.
\end{corollary}
\begin{proof}
    Проверим выполнение условия Маркова:
    \begin{multline*}
        {\D \overline{\xi_n}} = \dfrac{1}{n^2}{\sum\limits_{k = 1}^n \D\xi_k} \leq \\ \leq \dfrac{C}{n^2}\biggr(\sum\limits_{k = 1}^{n} k^\alpha\biggr) \leq \dfrac{C}{n^2}\int_0^{n} x^{\alpha} dx = \dfrac{C n^{\alpha + 1}}{n^2 (\alpha + 1)} = \\ = \dfrac{C}{n^{1 - \alpha}(\alpha + 1)} \ra 0, n \ra +\infty.
    \end{multline*}
    Первое равенство верно, так как случайные величины некоррелированные.
\end{proof}
\begin{reminder}
    Некоррелированность означает, что $\cov (\xi_i, \xi_j) = 0$ $\forall i \neq j$, и, вообще говоря, это не означает независимость!
\end{reminder}

\begin{theorem}[ЗБЧ в форме Хинчина]
    Пусть $\{\xi_n\}$ -- $i.i.d$: $\Ex \xi < +\infty$. Тогда выполнен ЗБЧ.
\end{theorem}
\begin{proof}
    Рассмотрим характеристическую функцию $\overline{\xi_n}$:
    \begin{multline*}
        \phi_{\overline{\xi_n}}(t) = \prod\limits_{k = 1}^{n} \phi_{\xi_k/n}(t) = (\phi_{\xi}(t/n))^n = \left(1 + i \Ex \xi \cfrac{t}{n} + o\left(\cfrac{1}{n}\right)\right)^n \ra e^{it \Ex \xi} = \phi_{\Ex \xi} (t), n \ra \infty \Longrightarrow \\
        \Longrightarrow \overline{\xi_n} \xrightarrow{d} \Ex \xi.
    \end{multline*}
    А в силу сходимости по распределению к константе мы получаем сходимость к ней и по вероятности, что и является выполнением ЗБЧ.
\end{proof}
\subsubsection{<<Метризация>> различных типов сходимости.}
\begin{theorem}[Метризуемость сходимости по вероятности.]
    Для последовательности случайных величин $\{Z_n\}_{n = 1}^{\infty}$ и случайной величины $Z$ следующие условия эквивалентны:
    \begin{enumerate}
        \item $Z_n \xrightarrow{\P} Z, n \ra \infty$.
        \item $d(Z_n, Z) = \Ex\biggr( \dfrac{|Z_n - Z|}{1 + |Z_n - Z|}\biggr) \ra 0, n \ra +\infty$.
    \end{enumerate}
    Причём $d(\cdot, \cdot)$ -- действительно является метрикой на пространстве случайных величин.
\end{theorem}
\begin{proof}
    Покажем $1) \Ra 2)$. \\
    Очевидно, что $\frac{|Z_n - Z|}{1 + |Z_n - Z|} \leq 1$ и что $\frac{|Z_n - Z|}{1 + |Z_n - Z|} \leq |Z_n - Z|$.
    Заметим, что
    \begin{multline*}
        d(Z_n, Z) = \\ = \int\limits_{|Z_n - Z| \leq \epsilon} \dfrac{|Z_n - Z|}{1 + |Z_n - Z|}d\P + \int\limits_{|Z_n - Z| > \epsilon} \dfrac{|Z_n - Z|}{1 + |Z_n - Z|}d\P \leq \\ \leq \epsilon + \P(|Z_n - Z| > \epsilon) \ra 0, n \ra +\infty, \epsilon \ra +0.
    \end{multline*}
    Теперь покажем в обратную сторону. \\
    Докажем от противного. Пусть нет сходимости по мере. Тогда $\exists \epsilon > 0$ и $\delta > 0$ такие, что $\forall N \in \N \ \exists n \geq N$ такой, что $\P(|\xi_n - \xi| > \epsilon) > \delta$.
    В силу монотонного возрастания $g(y) = \frac{y}{1 + y}$, мы получаем
    \[
        d(Z_n, Z) \geq \int_{|Z_n - Z| > \epsilon} \dfrac{|Z_n - Z|}{1 + |Z_n - Z|}d\P \geq \dfrac{\delta \epsilon}{1 + \epsilon} > 0.
    \]
    Что противоречит предположению о $d(Z_n, Z) \ra 0, n \ra \infty, n \ra \infty$.
    Значит $Z_n \xrightarrow{\P} Z$.
\end{proof}
\begin{note}
    Нетрудно показать, что $d(\cdot, \cdot)$ -- действительно определяет метрику: симметричность, неотрицательность и эквивалентность равенства нулю совпадению случайных величин очевидны, а корректность неравенства треугольника оставляется читателю в качестве несложного упражения. \\
    Аналогичным образом мы можем метризовать сходимость в среднем; в ней в качестве метрики берётся метрика, порождённая $L_p$-нормой.
\end{note}
\begin{note}
    Также можно показать, что модуль мы можем заменить на произвольную положительную степень $\alpha$, и это также будет эквивалентно сходимости по вероятности (но, возможно, уже не будет метрикой). \\
    Покажем это.
    Заметим, что $Z_n \xrightarrow{\P} Z, n \ra \infty \Longleftrightarrow |Z_n - Z| \xrightarrow{\P} 0, n \ra \infty$.
    В силу свойства наследования сходимости $|Z_n - Z|^\alpha \xrightarrow{\P} 0, n \ra \infty$.
    А значит и $d(|Z_n - Z|^\alpha, 0) \ra 0, n \ra +\infty$ -- что нам и требовалось показать. \\
    Покажем теперь обратную импликацию. Поскольку из $L_1$ сходимости следует $\P$, то
    \[
        \dfrac{|Z_n - Z|^\alpha}{1 + |Z_n - Z|^\alpha} \xrightarrow{\P} 0, n \ra \infty.
    \]
    Заметим, что функция $g(y) = \frac{y}{1 + y}$ -- обратима и обратная функция также непрерывна: $y = \frac{g(y)}{1 - g(y)}$.
    Используем свойство наследования сходимости и получаем что $|Z_n - Z|^\alpha \xrightarrow{\P} 0, n \ra \infty \Ra |Z_n - Z|\xrightarrow{\P} 0, n \ra \infty$, что и требовалось показать.
\end{note}
\begin{note}
    Для слабой сходимости случайных величин с непрерывными функциями распределения определена метрика Леви -- просто как $\sup$-метрика на функциях распределения:
    \[
        \xi_n \xrightarrow{d} \xi, n \ra \infty \Longleftrightarrow d_L(\xi_n, \xi) = \sup\limits_{x \in \R} |F_{\xi_n}(x) - F_{\xi}(x)| \ra 0, n \ra +\infty.
    \]
\end{note}
\subsubsection{Критерий Колмогорова.}
\begin{theorem}[Колмогоров, критерий выполнимости ЗБЧ]
    Для последовательности случайных величин $\xi_n$ следующие условия эквивалентны:
    \begin{enumerate}
        \item Для $\{\xi_n\}$ выполнен ЗБЧ.
        \item $\Ex \left( \dfrac{(\overline{\overset{\circ}{\xi_n}})^2}{1 + (\overline{\overset{\circ}{\xi_n}})^2}\right) \ra 0, n \ra +\infty$.
    \end{enumerate}
    Если же $\xi_n$ являются, вдобавок ко всему, независимыми случайными величинами, то второе условие упрощается до
    \[
        \sum\limits_{k = 1}^{n} \Ex \biggr(\dfrac{ (\overset{\circ}{\xi_k})^2}{1 + (\overset{\circ}{\xi_k})^2}\biggr) \ra 0, n \ra +\infty.
    \]
\end{theorem}
\begin{proof}
    Доказательство заключается в применении примечания к метризации сходимости по вероятности при $\alpha = 2$. \\
    Случай независимых случайных величин мы получаем из того, что применение борелевских функций не портит независимость случайных величин.
\end{proof}
\subsubsection{Неравенство Колмогорова.}
\begin{lemma}[Лемма Колмогорова]
    Пусть дана последовательность независимых случайных величин $\{\xi_n\}$.
    Обозначим за $S_n$ частичную сумму первых $n$ центрированных случайных величин последовательности.
    Справедливо следующее неравенство:
    \[
        \P\biggr(\sup\limits_{k = \overline{1, n}}|S_k| > \epsilon\biggr) \leq \dfrac{\Ex S_n^2}{\epsilon^2} = \dfrac{\D S_n}{\epsilon^2} = \dfrac{\sum_{k = 1}^{n} \D \xi_k}{\epsilon^2}.
    \]
\end{lemma}
\begin{proof}
    Зафиксируем некоторое $\epsilon > 0$.
    Введём следующую последовательность множеств $A_k$:
    \[
        A_k = \bigcap\limits_{j = 1}^{k - 1}\{w \mid |S_j(w)| \leq \epsilon\} \cap \{w \mid|S_{k}(w)| > \epsilon \}.
    \]
    Очевидно, что $i \neq j \Ra A_{i}A_j = \emptyset$.
    Пусть
    \[
        \bigcup\limits_{k = 1}^n A_k = A = \{w \in \Omega \mid \sup\limits_{k = \overline{1, n}} |S_k| > \epsilon\}.
    \]
    Заметим, что для индикаторов $A_k$ (ввиду их дизъюнктивности) верно следующее:
    \[
        I_{A} = \sum\limits_{k = 1}^n I_{A_k}.
    \]
    Так как
    \[
        \Ex S_n^2 = \int_{\Omega} S^2_n d\P \geq \int_{A} S^2_n d\P = \sum\limits_{k = 1}^{n} \int_{A_k} S^2_n d\P.
    \]
    Заметим, что для любого $k$ выполняется
    \[
        \Ex S_n^2 I_{A_k} \geq \Ex S_{k}^2 I_{A_k}.
    \]
    И действительно:
    \[
        \Ex S_n^2 I_{A_k} = \Ex S_k^2 I_{A_k} + \Ex \sum\limits_{i, j = k + 1}^{n} \Ex \biggr(\overset{\circ}{\xi_i} \overset{\circ}{\xi_j} I_{A_k}\biggr) + \Ex \biggr( S_k \sum\limits_{i = k + 1}^{n} \overset{\circ}{\xi_i} I_{A_k}\biggr) \geq \Ex S_k^2 I_{A_k}.
    \]
    (переход сделан в силу независимости $S_k$ от $\xi_{i}$ при $i > k$ и независимости $I_{A_k}$ от $\xi_{i}$, аналогично, при $i \geq k + 1$).
    Но на $A_k$ верно, что $S_k^2 > \epsilon^2$, а значит
    \[
        \Ex S_n^2 = \sum\limits_{k = 1}^{n} \Ex S_n^2 I_{A_k} \geq \sum\limits_{k = 1}^{n} \Ex S_k^2 I_{A_k} \geq \sum\limits_{k = 1}^{n} \P(A_k)\epsilon^2 = \P(A) \epsilon^2.
    \]
    На этом доказательство леммы завершено.
\end{proof}
\subsubsection{Теорема Колмогорова-Хинчина.}
\begin{theorem}[Колмогоров, Хинчин]
    Пусть дана последовательность независимых случайных величин $\{\xi_n\}$ таких, что $\sum\limits_{k = 1}^{+\infty} \D \xi_k < +\infty$.
    Введём следующее обозначение: $S_n = \sum\limits_{k = 1}^{n} \overset{\circ}{\xi_k}$.
    Тогда существует такая случайная величина $S$, что $S_n \xrightarrow{a.s.} S, n \ra \infty$.
\end{theorem}
\begin{proof}
    Зафиксируем $\epsilon > 0$.
    Определим последовательность множеств следующим образом:
    \[
        A^n_N = \{w \in \Omega \mid \max\limits_{k = 1, N} |S_{n + k} - S_n| > \epsilon\}.
    \]
    При всяком фиксированном $n$ эта последовательность не убывает по $N$.
    Объединение при фиксированном $n$ является следующим множеством:
    \[
        A^n = \bigcup\limits_{k = 1}^{+\infty} A^n_k = \{\sup\limits_{k \in \N} |S_{n + k} - S_n| > \epsilon\}.
    \]
    По теореме о непрерывности меры
    \[
        \P(A^n) = \lim \P(A^n_N), N \ra \infty.
    \]
    С другой стороны,
    \[
        \P(A^n_N) = \P\biggr(\max\limits_{k = 1, N} \biggr|\sum\limits_{j = 1}^{k} \overset{\circ}{\xi}_{j + n}\biggr| > \epsilon\biggr).
    \]
    Применим неравенство Колмогорова и получим
    \[
        \P(A^n_N) \leq \dfrac{1}{\epsilon^2}{\sum_{j = n}^{n + N} \D \xi_j} \ra \dfrac{1}{\epsilon^2}{\sum_{j = n}^{+\infty} \D \xi_j}, N \ra +\infty.
    \]
    Таким образом
    \[
        \P(A^n) \leq \dfrac{1}{\epsilon^2}{\sum_{j = n}^{+\infty} \D\xi_j}.
    \]
    Заметим, теперь, что последовательность $S_n$ -- фундаментальна почти наверное, поскольку остаток сходящегося ряда стремится к нулю.
    В силу критерия Коши сходимости почти наверное существует $S$ такая, что $S_n \xrightarrow{a.s.} S, n \ra \infty$.
\end{proof}
\subsubsection{Теоремы Колмогорова о двух и трёх рядах (доказательство достаточности).}
\begin{theorem}[Колмогоров, о двух рядах.]
    Пусть есть последовательность независимых случайных величин $\{\xi_n\}$ такая, что
    \[
        \sum\limits_{k = 1}^{+\infty} \Ex \xi_k < +\infty, \quad \sum\limits_{k = 1}^{+\infty} \D \xi_k < +\infty.
    \]
    Тогда $\sum\limits_{k = 1}^{+\infty} \xi_k < \infty$ (почти наверное).
\end{theorem}
\begin{proof}
    По теореме Колмогорова-Хинчина ряд $\sum\limits_{k = 1}^{+\infty} \xi_k - \Ex \xi_k$ -- сходится почти наверное.
    В силу сходимости ряда из математических ожиданий величин последовательности, почти наверное сходится и ряд из случайных величин.
\end{proof}
\begin{theorem}[Колмогоров, о трёх рядах.]
    Пусть есть последовательность независимых случайных величин $\{\xi_n\}$ и $\exists c > 0$ такое, что для последовательности усечённых случайных величин $\xi^c_n$, где срезка определяется как
    \[
        \xi_n^c(w) = \begin{cases}
                      \xi_n(w), & |\xi_n(w)| < c \\
                         0, & |\xi_n(w)| \geq c.
        \end{cases}
    \]
    выполнены условия теоремы Колмогорова о двух рядах:
    \[
        \sum\limits_{k = 1}^{+\infty} \Ex \xi_k^c < +\infty, \quad \sum\limits_{k = 1}^{+\infty} \D \xi_k^c < +\infty.
    \]
    И сходится ещё один ряд:
    \[
        \sum\limits_{k = 1}^{+\infty} \P(|\xi_n| \geq c) < +\infty.
    \]
    Тогда $\sum\limits_{k = 1}^{+\infty} \xi_k$ сходится почти наверное.
\end{theorem}
\begin{proof}
    По теореме Колмогорова о двух рядах ряд $\sum\limits_{k = 1}^{+\infty} \xi_k^c$ сходится почти наверное. \\
    С другой стороны,
    \[
        |\xi_n| \geq c \Longleftrightarrow \xi_n \neq \xi_n^c.
    \]
    По лемме Бореля-Кантелли
    \[
        \P(\limsup \{|\xi_n| \geq c\}) = 0.
    \]
    То есть с вероятностью единица
    \[
        \sum\limits_{k = 1}^{+\infty} I_{|\xi_k| \geq c} < +\infty.
    \]
    Тогда $\xi_n = \xi_n^c$ (почти наверное) для всех $n$ за исключением, быть может, конечного числа.
    А значит и ряд $\sum\limits_{k = 1}^{+\infty} \xi_k$ также сходится почти наверное.
\end{proof}
\begin{note}
    Для равномерно ограниченных последовательностей случайных величин мы можем обратить посылки теорем Колмогоров о двух и трёх рядах -- то есть если ряд сходится почти наверное, то в первом случае сойдутся ряды из математических ожиданий и дисперсий для последовательности, а во втором -- для всякого $c$ сойдутся ряды из математических ожиданий и дисперсий усечённых случайных величин.
\end{note}
\subsubsection{Усиленный закон больших чисел. Теоремы Колмогорова.}
Будем говорить, что для данной последовательности случайных величин $\xi_n$ выполнен усиленный закон больших чисел, если $\exists \{a_n\}$ такая, что
\[
    \overline{\xi_n} - a_n \xrightarrow{a.s.} 0, n \ra \infty.
\]
Традиционно, мы будем полагать $a_n = \Ex \overline{\xi_n}$, и тогда условием выполнимости УЗБЧ будет сходимость последовательности центрированных средних к нулю почти наверное.
\begin{lemma}[Теплиц]
    Пусть $\{a_n\}$ -- последовательность неотрицательных чисел, $b_n = \sum_{i = 1}^n a_i, b_1 = a_1 > 0$ и $b_n \ra \infty, n \ra \infty$. Пусть $x_n$ -- числовая последовательность, сходящаяся к $x \in \R$. Тогда
    \[
        \dfrac{1}{b_n}\sum\limits_{j = 1}^{n} a_j x_j \ra x, n \ra \infty.
    \]
\end{lemma}
\begin{proof}
    Зафиксируем $\epsilon > 0$. Поскольку $x_n \ra x$, то $\exists N = N(\epsilon)$ такое, что $\forall n \geq N \hookrightarrow$
    \[
        |x_n - x| \leq \dfrac{\epsilon}{2}.
    \]
    Поскольку $b_n \ra \infty, n \ra \infty$, то $\exists N^* > N$ такой, что 
    \[
        \dfrac{1}{N^*}\sum\limits_{j = 1}^N a_j|x_j - x| < \dfrac{\epsilon}{2}.
    \]
    Тогда $\forall n > N^* \hookrightarrow$
    \begin{multline*}
        \biggr|\dfrac{1}{b_n} \sum\limits_{j = 1}^n a_jx_j - x\biggr| \leq \dfrac{1}{b_n} \sum\limits_{j = 1}^n a_j|x_j - x| = \\ = \dfrac{1}{b_n}\sum\limits_{j = 1}^{N} a_j|x_j - x| + \dfrac{1}{b_n}\sum\limits_{j = N + 1}^n a_j|x_j - x| \leq \dfrac{1}{N^*}\sum\limits_{j = 1}^N a_j|x_j - x| + \dfrac{1}{b_n} a_j|x_j - x| \leq \\ \leq  \dfrac{\epsilon}{2} + \dfrac{b_n - b_N}{b_n}\dfrac{\epsilon}{2} \leq \epsilon.
    \end{multline*}
\end{proof}
\begin{lemma}[Кронекер]
    Пусть $\{b_n\}$ -- последовательность положительных возрастающих чисел такая, что $b_n \ra \infty, n \ra \infty$ и $\{x_n\}$ -- последовательность чисел таких, что $\sum_{j = 1}^\infty x_j$ сходится. Тогда
    \[
        \dfrac{1}{b_n}\sum\limits_{j = 1}^n b_j x_j \ra 0, n \ra \infty.
    \]
\end{lemma}
\begin{proof}
    Положим $b_0  = 0, S_0 = 0$ и $S_n = \sum_{j = 1}^n x_j$. Используем преобразование Абеля для $\sum_{j = 1}^n b_jx_j$:
    \[
        \sum\limits_{j = 1}^n b_j x_j = \sum\limits_{j = 1}^{n} b_j(S_j - S_{j - 1}) = b_nS_n - b_0S_0 - \sum\limits_{j = 1}^n S_{j - 1}(b_j - b_{j - 1}) = (*)
    \]
    Пусть $a_j = b_j - b_{j - 1}$. Тогда мы получаем
    \[
        (*) = S_n - \dfrac{1}{b_n}\sum\limits_{j = 1}^{n} S_{j - 1} a_j = (*)
    \]
    Так как $S_n \ra x, n \ra \infty$, то, по лемме Теплица,
    \[
        \dfrac{1}{b_n}\sum\limits_{j = 1}^{n} S_{j - 1}a_j \ra x, n \ra \infty.
    \]
    А значит
    \[
        (*) \ra x - x = 0, n \ra \infty.
    \]
\end{proof}
\begin{theorem}[Первая теорема Колмогорова]
    Если для последовательности независимых случайных величин $\{\xi_n\}$ ряд $\sum \dfrac{\D \xi_n}{n^2}$ сходится, то для неё выполнен УЗБЧ.
\end{theorem}
\begin{proof}
    Пусть $\eta_n = \frac{\xi_n - \Ex \xi_n}{n}$. Тогда, из условия теоремы, ряд $\sum\limits_{j = 1}^\infty \D \eta_j < +\infty$:
    \[
        \sum\limits_{j = 1}^{\infty} \D\eta_j =  \sum\limits_{j = 1}^{\infty} \dfrac{\D \xi_j}{j^2} < +\infty.
    \]
    По теореме Колмогорова о двух рядах $\sum\limits_{j = 1}^\infty \eta_j$ сходится почти наверное.
    Применим лемму Кронекера к $b_n = n$ и получим:
    \[
        \dfrac{1}{n}\sum\limits_{j = 1}^{n} \eta_j \cdot j \xrightarrow{a.s.} 0, n \ra \infty.
    \]
    Но это и означает выполнение УЗБЧ, поскольку $j\eta_j = \xi_j - \Ex \xi_j$:
    \[
        \dfrac{1}{n}\sum\limits_{j = 1}^{n} \xi_j - \Ex \xi_j \xrightarrow{a.s.} 0, n \ra \infty.
    \]
\end{proof}
\begin{theorem}[Вторая теорема Колмогорова.]
    Если $\xi_n$ -- $i.i.d$ последовательность, то следующие условия эквивалентны:
    \begin{enumerate}
        \item $\Ex \xi = a$.
        \item Для $\{\xi_n\}$ выполнен УЗБЧ с $a_n = a$.
    \end{enumerate}
\end{theorem}
\begin{proof}
    Пусть $\exists\  \Ex|\xi| < +\infty$. \\
    Рассмотрим последовательность $\{\xi_n^*\}$, случайные величины в которой определим как
    \[
        \xi_n^*(w) =
        \begin{cases}
            \xi_n(w), & |\xi_n(w)| < n \\
            0, & |\xi_n(w)| \geq n.
        \end{cases}
    \]
    Оценим дисперсию $\xi_n^*$:
    \begin{multline*}
        \D \xi_n^* \leq \Ex (\xi_n^*)^2 = \int_{|x| < n} x^2dF_\xi = \sum\limits_{k = 0}^{n - 1}\int_{k \leq |x| < k + 1}x^2dF_\xi \leq \\ \leq \sum\limits_{k = 0}^{n - 1}(k + 1)^2 \int_{k \leq |x| < k + 1}dF_{\xi} = \sum\limits_{k = 0}^{n - 1}(k + 1)^2\P(k \leq |\xi| < k + 1) = (*)
    \end{multline*}

    Введём следующее обозначение для сокращения записи: $\P(k \leq |\xi| < k + 1) = [k, k + 1)$. $(*)$ перепишется как
    \[
        (*) = \sum\limits_{k = 0}^{n - 1}(k + 1)^2[k, k + 1).
    \]
    Теперь исследуем ряд $\sum_{n = 1}^{\infty} \frac{1}{n^2}\D\xi_n^*$ на сходимость:
    \[
        \sum\limits_{n = 1}^{\infty} \dfrac{\D \xi_n^*}{n^2} \leq \sum\limits_{n = 1}^{\infty} \dfrac{1}{n^2}\sum\limits_{k = 0}^{n - 1}(k + 1)^2[k, k + 1) = (**).
    \]
    Заметим, что все члены ряда -- неотрицательны. Тогда перестановка членов ряда не влияет на сходимость и не изменяет сумму ряда. Поэтому мы можем представить $(**)$ в виде:
    \begin{multline*}
        (**) = \sum\limits_{k = 0}^{\infty} (k + 1)^2 [k, k + 1) \sum\limits_{n = k + 1}^{\infty} \dfrac{1}{n^2} \leq 3[0, 1) + \sum\limits_{k = 1}^{\infty}(k + 1)^2[k, k + 1) \sum\limits_{n = k + 1}^{\infty}\dfrac{1}{n^2} \leq \\ \leq 3[0, 1) + \sum\limits_{k = 1}^{\infty}(k + 1)^2 [k, k + 1)\int_k^{+\infty}\dfrac{1}{x^2}dx \leq 3[0, 1] + \sum\limits_{k = 1}^{\infty} \dfrac{(k + 1)^2}{k}[k, k + 1) \leq (**)
    \end{multline*}
    Заметим, что при $k \geq 1$ выполняется следующее неравенство: $(k + 1)^2 \leq k(k + 3)$, то есть 
    \[
        \dfrac{(k + 1)^2}{k} \leq k + 3.
    \]
    Поэтому $(**)$ оценивается сверху как 
    \begin{multline*}
        (**) \leq 3\sum\limits_{k = 0}^{\infty}[k, k + 1) + \sum\limits_{k = 1}^{\infty} k [k, k + 1) \leq  3 + \sum\limits_{k = 0}^{\infty}k \int_{k \leq |x| < k + 1}dF_{\xi} \leq \\ \leq 3 + \sum\limits_{k = 0}^{\infty}\int_{k \leq |x| < k + 1} |x|dF_{\xi} = 3 + \Ex|\xi| < +\infty.
    \end{multline*}
    А значит, по первой теореме Колмогорова, имеет место следующая сходимость: \[\overline{\xi_n^*} - \Ex \overline{\xi_n^*} \xrightarrow{a.s.} 0, n \ra \infty.\]
    Исследуем на сходимость следующий ряд:
    \[
        \sum\limits_{n = 1}^{\infty} \P(\xi_n^* \neq \xi_n).
    \]
    Если он сходится, то по лемме Бореля-Кантелли, ряды $\overline{\xi_n^*}$ и $\overline{\xi}$ отличаются лишь конечным числом членов, и сходятся или расходятся одновременно (где сходимость понимается в смысле сходимости почти наверное). \\
    По определению $\xi_n^*$ мы можем расписать $\P(\xi_n^* \neq \xi_n)$ как 
    \[
        \P(\xi_n^* \neq \xi_n) = \P(|\xi| \geq n) = \sum\limits_{k = n}^{\infty}[k, k + 1).
    \]
    Поэтому 
    \[
        \sum\limits_{n = 1}^{\infty}\P(\xi_n^* \neq \xi_n) = \sum\limits_{n = 1}^{\infty}\sum\limits_{k = n}^{\infty} [k, k + 1) = \sum\limits_{k = 1}^{\infty} k [k, k + 1) \leq \Ex|\xi| < +\infty.
    \]
    Значит имеет место следующая сходимость:
    \[
        \overline{\xi_n} - \Ex \overline{\xi_n^*} \xrightarrow{a.s.} 0, n \ra \infty.
    \]
    В тоже время, по теореме Хинчина, $\overline{\xi_n} \xrightarrow{\P} \Ex \xi, n \ra \infty$ и, в силу единственности предела, импликация $2) \Ra 1)$ показана:
    \[
        \overline{\xi_n} - \Ex \xi \xrightarrow{a.s.} 0, n \ra \infty.
    \]
    Пусть теперь $\overline{\xi_n} - a \xrightarrow{a.s.} 0, n \ra \infty$. \\
    Заметим, что тогда $\frac{1}{n}\xi_n \xrightarrow{a.s.}0, n \ra \infty$. 
    И действительно:
    \[
        \dfrac{\xi_n}{n} = \dfrac{1}{n}\sum\limits_{k = 1}^n \xi_k - \dfrac{n - 1}{n} \cdot \dfrac{1}{n - 1}\sum\limits_{k = 1}^{n - 1}\xi_k \xrightarrow{a.s.} a - 1 \cdot a = 0, n \ra \infty.
    \]
    А значит, поскольку события $\biggr\{|\frac{\xi_n}{n}| > 1\biggr\}$ -- независимы в совокупности, то ряд 
    \[
        \sum\limits_{n = 1}^{\infty} \P\biggr(\frac{\xi_n}{n} \geq 1\biggr) < +\infty.
    \]
    иначе по второй части леммы Бореля-Кантелли $\frac{\xi_n}{n}$ не может п.н. сходится к 0. \\
    С другой стороны
    \begin{multline*}
        +\infty > \P(|\xi| \geq 0) + \sum\limits_{n = 1}^{\infty} \P(|\xi_n| \geq n) = \sum\limits_{n = 0}^{\infty} \P(|\xi| \geq n) = \sum\limits_{k = 0}^{\infty} (k + 1)[k, k + 1) = \\ = \sum\limits_{k = 0}^{\infty}(k + 1)\int_{k \leq |x| < k + 1}dF_\xi \geq \sum\limits_{k = 0}^{\infty} \int_{k \leq |x| < k + 1} |x|dF_\xi = \Ex|\xi|.
    \end{multline*}
    А значит у $\xi$ существует математическое ожидание. Более того, тогда, по теореме Хинчина, $\overline{\xi_n} \xrightarrow{\P} \Ex \xi$ и $a = \Ex \xi$.
\end{proof}
